\documentclass[10pt,a4paper]{amsart}
\usepackage[margin=19mm]{geometry}
\usepackage{amsmath,amssymb,mathtools}
\usepackage[scr=rsfs,cal=euler]{mathalfa}
\usepackage{xfrac}
\usepackage[unicode,hidelinks,bookmarks=false]{hyperref}
\newcommand{\ie}{i.e.\ }
\newcommand{\cf}{cf.\ }
\newcommand{\resp}{respectively\ }
\newcommand{\Subsection}{\subsection}
\providecommand{\nobreakdash}{\nobreak\hbox{-}}
\setlength{\emergencystretch}{2em}
\multlinegap0pt
\usepackage{bm}
\usepackage{bbm}
\usepackage{dsfont}
\usepackage{xspace}
\usepackage{cancel}
\usepackage{mathtools}
\usepackage{amsthm}
\makeatletter
\@ifundefined{theorem}{\newtheorem{theorem}{Theorem}}{}
\@ifundefined{lemma}{\newtheorem{lemma}[theorem]{Lemma}}{}
\makeatother
\DeclareMathOperator{\M}{M}
\newcommand{\bQ}{\mathbb Q}
\title[Twisted Jacquet modules associated to maximal parabolic subg]{Twisted Jacquet modules associated to maximal parabolic subgroups and cuspidal representations of $GL(n, q)$}
\author{Kumar Balasubramanian, Krishna Kaipa, Himanshi Khurana}
\date{Source: arXiv:2606.22846v2 (CC BY 4.0)}
\begin{document}
\maketitle

\noindent\emph{Source and licence.} Twisted Jacquet modules associated to maximal parabolic subgroups and cuspidal representations of $GL(n, q)$. Version arXiv:2606.22846v2 (CC BY 4.0), distributed under the Creative Commons Attribution 4.0 International licence, as recorded in the arXiv licence metadata for this exact version and shown on the abstract page https://arxiv.org/abs/2606.22846v2. Full rights notice: licenses/arxiv-2606.22846-CC-BY-4.0.txt.\par\smallskip

\noindent\textit{Editorial scope.} Counted unit: the Lemma whose source label is \texttt{amnr\_lemma} in the article (source0.tex lines 489--492, source label \texttt{amnr\_lemma}), delivered together with the article's own complete original proof of it (source0.tex lines 493--522, one \texttt{proof} environment of 30 lines). No mathematics is written, repaired or re-derived: statement and proof are the article's own text line for line. Required context retained uncounted: none. Every definition, display and label of those lines is retained unchanged. Omitted from this package and not redistributed: 1--488, 523--1344. Nothing inside a retained excerpt is omitted or altered: every retained body is delivered exactly as the article's lines are written, so no omission receipt of any kind accompanies this package. Citations: the retained excerpts carry no citation command at all, so no bibliography entry of the article is transcribed and no reference is redistributed. Reference adaptation: none was needed and none was made; every reference inside the delivered unit resolves inside it, so no unresolved reference or citation is produced. Printed slips kept literal: no printed word, symbol, hypothesis or number of the counted statement or of its proof was altered. Layout and class: the article's own class is \texttt{amsart} with packages including babel, url, graphicx, amsopn, amsmath, amsfonts, amssymb, tikz-cd, amsthm, xfrac, mathdots, setspace, mathtools, multirow; the wrapper is the lane's pinned \texttt{amsart} editorial wrapper at 10pt on A4, which restates the theorem-like environments the retained text needs and adds the article's own macro definitions that the retained text needs (\texttt{\textbackslash M}, \texttt{\textbackslash bQ}), with extra wrapper packages (bm, bbm, dsfont, xspace, cancel, mathtools). The delivered numbering is the unit's own. Assets: no retained span contains a figure environment, a graphics-inclusion command, a PDF-page inclusion, a table or a TikZ picture; no image file is copied into this package, the package declares no asset dependency and no image file is redistributed with it. Licence: version arXiv:2606.22846v2 of this article is distributed under the Creative Commons Attribution 4.0 International licence, as recorded in the arXiv licence metadata for this exact version and shown on the abstract page https://arxiv.org/abs/2606.22846v2. A full rights notice is delivered at licenses/arxiv-2606.22846-CC-BY-4.0.txt. Characters: every retained excerpt is pure ASCII, so the delivered engine, XeLaTeX with its default Latin Modern text font, reports no missing character.
\begin{lemma} \label{amnr_lemma}
The following polynomial identity holds in the polynomial ring  $\bQ[S]$:
\[\sum_r a(m \times n,r, q) (S;q)_{n-r}= q^{mn} (Sq^{-m};q)_n.\] 
\end{lemma}

\begin{proof}
For $\ell \geq n$, the size of the set $\M_{n \times \ell}(n,F)$ of all full rank $n \times \ell $ matrices is 
 \[ a( n\times \ell, n,q)=\prod_{i=0}^{n-1} (q^\ell-q^i).\]
 For such a matrix $X \in \M_{n \times \ell}(n,F)$, and for $m \leq \ell$,  let $X=\begin{bmatrix}  X_1 &  X_2 \end{bmatrix}$ where $X_1$ has size $n \times m$ and $X_2$ has size $ n \times (\ell-m)$. 
For $X_1$ of rank $r$, let $U$ denote the column span of $X_1$. The number of matrices $X_2$ such that 
$\begin{bmatrix}  X_1 &  X_2 \end{bmatrix}$ has rank $n$, is same as the number of maps $F^{\ell-m} \to F^n$ such that the associated map $F^{\ell-m} \to F^n/U$ has rank $(n-r)$. In other words, the number of choices for $X_2$ is \[ q^{(\ell-m)r} \cdot a(n-r \times \ell-m, n-r, q)= q^{(\ell-m)r} \cdot \prod_{i=0}^{n-r-1}(q^{\ell-m}-q^i). \]
 
 
 Therefore, we get the identity
\[  \prod_{i=0}^{n-1} (q^\ell-q^i)=\sum_{r} a(m \times n, r, q) q^{(\ell-m)r} \prod_{i=0}^{n-r-1}(q^{\ell-m}-q^i). \]
Multiplying 
 by $q^{mn}$ we can rewrite this as 
 \[  q^{mn} \prod_{i=0}^{n-1} (q^\ell-q^i)=\sum_{r} a(m \times n, r, q) q^{\ell r} \prod_{i=0}^{n-r-1}(q^{\ell}-q^{i+m}). \]
This implies the polynomial identity in $\bQ[T]$:
\[
q^{mn} \prod_{i=0}^{n-1} (T-q^i)=
\sum_{r} a(m \times n, r, q) T^{r} \prod_{i=0}^{n-r-1}(T-q^{i+m}),
\]
because the difference between the RHS and LHS is a polynomial of degree at most $n$ in $\bQ[T]$ with infinitely many roots $T=q^\ell$ for $\ell \geq n$, and hence it is the  zero polynomial. Dividing by $T^n$, we get the identity in $\bQ[T^{-1}]$:
\[
q^{mn} \prod_{i=0}^{n-1} (1 -q^i T^{-1})=
\sum_{r} a(m \times n, r, q)   \prod_{i=0}^{n-r-1}(1-q^{i+m}T^{-1}).
\]
In terms of $S=q^{m}T^{-1}$ we can rewrite this as the identity in $\bQ[S]$ given by:
\[
q^{mn}(Sq^{-m};q)_n=
\sum_{r} a(m \times n, r, q) \, 
(S;q)_{n-r}.
\]
\end{proof}

\end{document}
